\chapter{深度生成模型}
\begin{quote}\small
\textit{本章摘要。}本章围绕一个问题展开：如何让模型生成符合数据规律、又能用于具体任务的新样本？先区分采样、补全、去噪和条件生成，再以变分自编码器（VAE）、生成对抗网络（GAN）和扩散模型为里程碑，说明它们的生成机制、训练目标和典型失效方式。最后讨论任务选择、样本验证和工业使用边界。
\end{quote}

生成模型不能替真实观测，也不能凭空证明罕见故障确实存在。它可以在数据稀缺时产生候选样本，补全缺失观测，或在给定工况下模拟可能的退化轨迹。合成数据是否有用，要在未参与训练和模型选择的真实数据上检验。

\section{生成任务与概率建模目标}

\subsection{先确定要生成什么}

假设$x\in\R^D$是一个振动窗口，$c$是转速、负载或故障类别等条件。深度生成模型的常见任务及相关用途包括：
\begin{enumerate}
  \item \textbf{无条件采样：}从学习到的分布中产生新的$x$，例如模拟正常振动窗口。
  \item \textbf{条件生成：}给定$c$产生$x$，例如在指定转速下生成可能的波形。条件通常允许多个结果，不要求给定条件后输出唯一样本。
  \item \textbf{补全与去噪：}给定部分观测、带噪观测或掩码，生成与已知信息一致的完整结果。
  \item \textbf{表示学习：}把高维样本压缩为较低维的潜变量，供可视化、聚类或下游预测使用。
\end{enumerate}

这些任务的共同对象是分布。若模型学习$p_{\theta}(x)$，就可以从该分布采样；若学习$p_{\theta}(x\mid c)$，就可以在条件$c$下采样。训练集是有限观测，不等于真实分布$p_{\mathrm{data}}$。连续变量的$p_{\theta}(x)$是密度，某个区域的概率需要对密度积分。把所有样本都替换成均值，只能描述平均位置，不能再现波动、相关性和少数模式。

这里$\theta$表示模型的可训练参数。固定观测$x$、把$p_{\theta}(x)$看作参数的函数，就得到\textbf{似然}。一种常见训练原则是最大似然，即增大训练样本的平均对数似然；但不同生成模型实现这一目标的方式不同，也有模型采用其他训练准则。训练时利用数据更新参数；采样时固定参数，利用随机输入产生新样本。

生成目标与判别目标不同。判别模型通常学习$p(y\mid x)$或直接学习决策边界，生成模型还要描述$x$本身的变化。生成模型的训练损失也不自动等于任务评价指标。重构误差较小，不代表裂纹位置正确；样本看起来逼真，也不代表它覆盖了真实故障类型。

本章聚焦 VAE、GAN 和扩散模型。它们并不涵盖所有生成模型：自回归模型按顺序预测下一个分量的条件分布，正规化流通过可逆变换及变量替换公式计算密度。先掌握本章的三条主线，有助于理解不同模型如何安排训练目标和采样过程。

\subsection{潜变量把随机性变成可采样结构}

潜变量$z$是模型引入、但数据中没有直接观测的变量。一个常见生成方向是
\[
z\sim p(z),\qquad x\sim p_{\theta}(x\mid z).
\]
先抽取简单先验中的$z$，再由解码器把它变成样本$x$。例如，若$z\sim\mathcal N(0,1)$且$x\mid z\sim\mathcal N((z,2z)^{\mathsf T},\sigma_x^2I)$，两个观测分量会因共享$z$而相关。观察到$x$后反推$z$，对应后验$p_{\theta}(z\mid x)$。
这里的先验$p(z)$是观察$x$之前对潜变量的分布设定，后验则描述观察$x$之后对$z$的判断。下文通常令$z\in\R^d$；$I$表示上下文相应维度的单位矩阵，$I_d$则明确表示$d\times d$单位矩阵，$\mathcal N(\mu,\Sigma)$表示均值为$\mu$、协方差为$\Sigma$的高斯分布。

普通自编码器只要求编码器和解码器互相重构。它没有规定新样本应该从编码空间的哪里抽取。变分自编码器把潜变量先验、概率解码器和近似后验放在同一个目标中，从而得到可采样的表示。

\section{VAE：用潜变量获得可采样表示}

\subsection{为什么需要变分近似}

变分自编码器（Variational Autoencoder，VAE）用潜变量模型表示联合分布：
\begin{equation}
p_{\theta}(x,z)=p(z)p_{\theta}(x\mid z).
\end{equation}
给定$x$后，真实后验满足
\[
p_{\theta}(z\mid x)=\frac{p_{\theta}(x\mid z)p(z)}{p_{\theta}(x)},\qquad
p_{\theta}(x)=\int p_{\theta}(x\mid z)p(z)\,\mathrm dz.
\]
分母通常无法解析计算。VAE 用编码器给出的$q_{\phi}(z\mid x)$近似该后验，并最大化证据下界（Evidence Lower Bound，ELBO）：
\begin{equation}
\mathcal L(x)=\E_{q_{\phi}(z\mid x)}[\log p_{\theta}(x\mid z)]-\KL(q_{\phi}(z\mid x)\|p(z)).
\label{eq:generative-vae-elbo}
\end{equation}
这里$\phi$是编码器参数，$\E_q$表示按分布$q$取平均。第一项鼓励从近似后验抽到的$z$能够解释输入$x$；第二项用 KL 散度约束近似后验与先验的差异，使训练时使用的潜变量与生成时的先验采样相协调。两项需要共同优化，不能只把 KL 压到零。

由 Jensen 不等式，
\begin{align}
\log p_{\theta}(x)
&=\log\int q_{\phi}(z\mid x)\frac{p_{\theta}(x,z)}{q_{\phi}(z\mid x)}\,\mathrm dz\\
&\geq \E_{q_{\phi}}[\log p_{\theta}(x,z)-\log q_{\phi}(z\mid x)]
=\mathcal L(x).
\end{align}
这里假设相关积分和期望存在，且$q_{\phi}(z\mid x)$在$p_{\theta}(x,z)>0$的区域为正；下文采用的非退化高斯近似满足这一支持条件。
更精确的关系是
\begin{equation}
\log p_{\theta}(x)-\mathcal L(x)=\KL\bigl(q_{\phi}(z\mid x)\|p_{\theta}(z\mid x)\bigr)\geq0.
\end{equation}
这里的两个 KL 项用途不同。式\eqref{eq:generative-vae-elbo}中的 KL 比较近似后验与先验，最后一个等式中的 KL 衡量后验近似误差。KL 散度非负但通常不对称，也不是欧氏距离。
对连续分布，它的定义是
\[
\KL(q\|p)=\int q(z)\log\frac{q(z)}{p(z)}\,\mathrm dz.
\]
本章的$\log$均为自然对数；若$p$在$q$具有正概率的区域为零，KL 散度可以为无穷大。

\subsection{如何训练和采样}

常用的近似后验为
\[
q_{\phi}(z\mid x)=\mathcal N\bigl(\mu_{\phi}(x),\diag(\sigma_{\phi}^2(x))\bigr).
\]
其中$\mu_{\phi}(x),\sigma_{\phi}(x)\in\R^d$，标准差各分量为正，$\diag$把向量放到矩阵对角线上。若把抽样当作不透明操作，就不能直接沿抽样路径对编码器参数反向传播。重参数化把随机性移到与参数无关的$\epsilon$：
\begin{equation}
z=\mu_{\phi}(x)+\sigma_{\phi}(x)\odot\epsilon,
\qquad \epsilon\sim\mathcal N(0,I_d).
\end{equation}
符号$\odot$表示逐元素相乘。
因此固定一次$\epsilon$时，梯度可以通过$z$传回编码器。训练阶段给定$x$，从$q_{\phi}(z\mid x)$抽样并计算重构项。生成阶段没有$x$，改为从$p(z)$抽样，再交给解码器$p_{\theta}(x\mid z)$。

若解码器采用$p_{\theta}(x\mid z)=\mathcal N(\mu_{\theta}(z),\sigma_x^2I)$，且$\sigma_x^2>0$固定，则
\[
-\log p_{\theta}(x\mid z)=\frac{1}{2\sigma_x^2}\|x-\mu_{\theta}(z)\|_2^2+\frac D2\log(2\pi\sigma_x^2).
\]
因此，在这一概率解码器下，最小化负对数似然等价于最小化平方重构误差，除以$D$则得到均方误差。这里假定给定$z$后各观测分量独立且方差相同。这种设定对少数但重要的冲击或裂纹像素未必足够合适。仅输出$\mu_{\theta}(z)$是给定$z$后的均值预测；从上述高斯分布继续抽样，才是完整的概率解码过程。

当$q_{\phi}$是对角高斯、$p(z)=\mathcal N(0,I)$时，KL 项为
\[
\KL(q\|p)=\frac12\sum_{j=1}^{d}(\mu_j^2+\sigma_j^2-\log\sigma_j^2-1).
\]
例如一维情况下$\mu=1,\sigma=1$时，KL 为$0.5$。$\mu=0,\sigma=1$时 KL 为零，但这只说明两个分布相同，不说明重构准确。

\subsection{VAE 的变体与典型问题}

VAE 的主要折中是重构准确性与潜变量规整性。若对不同输入都有$q_{\phi}(z\mid x)\approx p(z)$，潜变量就几乎不再携带输入信息，这称为\textbf{后验坍缩}。过强的 KL 惩罚、能够忽略$z$的强解码器以及优化过程，都可能促成这一现象。KL 退火（KL annealing）在训练初期使用较小的 KL 权重，再逐渐增加；free bits 为潜变量各维或各组设置少量 KL 额度，在额度内不继续奖励减小 KL。这些方法改变优化折中，却不保证避免坍缩，也不保证潜变量对应可解释的物理因素。

条件 VAE 将条件加入编码器和解码器，使用$q_{\phi}(z\mid x,c)$和$p_{\theta}(x\mid z,c)$；潜变量先验可以仍取$p(z)$，也可以设为$p(z\mid c)$。生成时给定$c$，从相应先验抽取$z$，再解码得到$x$。它适合需要连续潜在表示和条件采样的任务。采用简单高斯解码器并输出其均值时，VAE 容易产生较平滑的结果：平方损失偏向给定潜变量后的条件均值，未被潜变量区分的细节可能被平均。若细节和少数模式重要，应检查模式覆盖，而不能只增加潜变量维度。

\section{GAN：追求逼真样本的对抗训练}

\subsection{判别器如何提供训练信号}

生成对抗网络（Generative Adversarial Network，GAN）把采样器和判别器分开。生成器$G:\R^d\to\R^D$把噪声$z$变成$G(z)$，判别器$D_{\psi}:\R^D\to(0,1)$估计输入来自真实数据的概率，其中$\psi$是判别器参数。训练时交替更新：判别器区分真实样本与合成样本，生成器则使合成样本更难被判别。采样时只需要生成器。

原始 GAN 目标为
\begin{equation}
\min_G\max_{D_{\psi}}\;\E_{x\sim p_{\mathrm{data}}}\log D_{\psi}(x)+\E_{z\sim p(z)}\log(1-D_{\psi}(G(z))).
\label{eq:generative-gan}
\end{equation}
固定生成分布$p_g$时，在$p_{\mathrm{data}}(x)+p_g(x)>0$处，逐点最优判别器为
\[
D^*(x)=\frac{p_{\mathrm{data}}(x)}{p_{\mathrm{data}}(x)+p_g(x)}.
\]
这里用$p_{\mathrm{data}}$和$p_g$表示相对于同一参考测度的密度；若最优值为$0$或$1$，可按判别器输出的极限理解。代回式\eqref{eq:generative-gan}，最优判别器下的目标值为
\[
C(G)=-\log4+2\operatorname{JSD}(p_{\mathrm{data}}\|p_g).
\]
其中 Jensen--Shannon 散度定义为
\[
\operatorname{JSD}(p\|q)=\tfrac12\KL(p\|m)+\tfrac12\KL(q\|m),
\qquad m=\tfrac12(p+q).
\]
它在两个分布相同时为零，因此理论最优点满足$p_g=p_{\mathrm{data}}$。这一结论假设对每个$G$都能取得最优判别器，并且生成器能够表示目标分布；它不保证有限网络的交替训练能达到该最优点。

\subsection{逼真与覆盖为何会冲突}

GAN 的训练和采样不需要显式计算$p_g(x)$，判别器提供学习数据结构的训练信号。它的困难来自两个网络的交替博弈。在原始极小极大目标中，若判别器很容易识别合成样本，生成器的梯度可能很弱。实践中常改为最小化非饱和生成器损失
\[
\mathcal L_G^{\mathrm{ns}}=-\E_{z\sim p(z)}\log D_{\psi}(G(z)),
\]
以改善这种训练信号不足的问题，但这并不消除所有训练困难。不同噪声被映射到少数相似样本、使生成分布遗漏真实模式时，称为\textbf{模式坍塌}。
生成结果可能很逼真，却遗漏罕见故障或某些工况。条件 GAN 可以把类别、工况或部分观测输入生成器和判别器，但条件一致性仍需单独验证。

GAN 适合采样速度要求高、且有足够样本支持目标分布的任务。它不提供像 VAE 那样直接的后验编码，也不直接给出可比较的显式似然。训练稳定性、模式覆盖和条件满足率应与样本质量一起报告。

\section{扩散模型：通过逐步去噪生成样本}

\subsection{从一次生成改为多步恢复}

扩散模型把困难的生成映射分成许多局部去噪任务。前向过程逐步破坏真实样本，反向网络学习在每个噪声尺度恢复结构。设$x_0\sim p_{\mathrm{data}}$，$t=1,\ldots,T$，并令$0<\beta_t<1$：
\begin{equation}
q(x_t\mid x_{t-1})=\mathcal N(\sqrt{1-\beta_t}x_{t-1},\beta_tI).
\end{equation}
记$\alpha_t=1-\beta_t$、$\bar\alpha_t=\prod_{s=1}^t\alpha_s$，可直接得到
\begin{equation}
x_t=\sqrt{\bar\alpha_t}x_0+\sqrt{1-\bar\alpha_t}\epsilon,
\qquad \epsilon\sim\mathcal N(0,I).
\label{eq:generative-direct-noising}
\end{equation}
这里$t$是人为设定的噪声步，不是设备运行时间。训练时随机选$t$和$\epsilon$，所以不必真的执行完整的前向链。
通常预先设定噪声日程$\{\beta_t\}_{t=1}^T$，使$\bar\alpha_T$足够接近零，此时终点分布才近似标准高斯。上式给出了单个噪声步的边缘采样方式；分别用独立$\epsilon$计算多个$x_t$，并不等于抽取一条具有正确步间相关性的前向轨迹。这里的$q$是固定加噪分布，与 VAE 中可学习的$q_{\phi}$用途不同。

例如$x_0=2$、$\bar\alpha_t=0.64$且$\epsilon=1$时，$x_t=0.8\times2+0.6\times1=2.2$。网络看到$x_t$和$t$，不知道$x_0$。它要学习的是在给定噪声状态下最有用的恢复方向。

\subsection{噪声预测与反向采样}

噪声预测网络记为$\epsilon_{\theta}(x_t,t)$，常用训练目标为
\begin{equation}
\mathcal L_{\mathrm{simple}}=\E_{t,x_0,\epsilon}\left[\|\epsilon-\epsilon_{\theta}(x_t,t)\|_2^2\right].
\end{equation}
这里通常均匀抽取$t\in\{1,\ldots,T\}$，从训练集中抽取$x_0$，再独立抽取$\epsilon\sim\mathcal N(0,I)$，用式\eqref{eq:generative-direct-noising}构造网络输入。训练程序知道抽取的噪声，并将它作为监督目标；网络的输入只有$x_t$和$t$。
平方损失的最优函数是条件均值：
\[
\epsilon_{\theta}^*(x_t,t)=\E[\epsilon\mid x_t,t].
\]
由式\eqref{eq:generative-direct-noising}，可以得到干净样本估计
\[
\widehat x_0=\frac{x_t-\sqrt{1-\bar\alpha_t}\epsilon_{\theta}(x_t,t)}{\sqrt{\bar\alpha_t}}.
\]
当$\bar\alpha_t$很小时，上式中的噪声预测误差会被放大，因此高噪声下的一次干净样本估计可能不准确。完整反向过程会在多个噪声尺度上继续修正，不能只用这一估计判断最终采样质量。

采样从$x_T\sim\mathcal N(0,I)$开始，再按$t=T,T-1,\ldots,1$反向计算。以去噪扩散概率模型（Denoising Diffusion Probabilistic Model，DDPM）的一种固定方差参数化为例，令$\bar\alpha_0=1$，并取
\begin{align}
\mu_{\theta}(x_t,t)&=\frac1{\sqrt{\alpha_t}}
\left(x_t-\frac{\beta_t}{\sqrt{1-\bar\alpha_t}}\epsilon_{\theta}(x_t,t)\right),\\
\widetilde\beta_t&=\frac{1-\bar\alpha_{t-1}}{1-\bar\alpha_t}\beta_t,\\
x_{t-1}&=\mu_{\theta}(x_t,t)+\sqrt{\widetilde\beta_t}\,\zeta_t,
\qquad t=T,\ldots,2.
\end{align}
其中$\zeta_t\sim\mathcal N(0,I)$是各反向步骤独立抽取的噪声，$\widetilde\beta_t$来自可解析的前向后验$q(x_{t-1}\mid x_t,x_0)$。在这里展示的采样约定中，最后一步直接输出$x_0=\mu_{\theta}(x_1,1)$，不再添加噪声。反向方差也可以采用其他设定。初始噪声已经提供随机性，逐步添加噪声并非产生多样样本的必要条件；一次计算$\widehat x_0$也不等于执行完整反向采样。

完整的变分目标包含前向后验与反向转移之间的 KL 项。常用的噪声均方误差是经过参数化和权重处理后的简化目标，不能与完整变分目标逐项等同。
除了预测噪声，也可以预测$x_0$，或预测二者的线性组合，例如$v=\sqrt{\bar\alpha_t}\epsilon-\sqrt{1-\bar\alpha_t}x_0$。这种$v$预测也称速度预测，但$v$并不是设备的物理速度。不同参数化与损失权重会影响训练和采样表现，但共享逐步去噪框架。

\subsection{条件控制与扩散模型的变体}

给定条件$c$时，网络写成$\epsilon_{\theta}(x_t,t,c)$。无分类器引导（Classifier-Free Guidance，CFG）通常在训练时随机丢弃一部分条件，让同一网络学习有条件和无条件预测；采样时将二者组合为
\begin{equation}
\widehat\epsilon=\epsilon_{\theta}(x_t,t,\varnothing)+w\left[\epsilon_{\theta}(x_t,t,c)-\epsilon_{\theta}(x_t,t,\varnothing)\right].
\end{equation}
这里$\varnothing$表示空条件。按本式约定，$w=0$是无条件预测，$w=1$是普通条件预测，$w>1$会放大两者之差。较强引导通常增强条件影响，但可能降低多样性并产生伪影，也不再一般对应原条件分布的精确采样。
边缘、深度、掩码或几何信息也可以作为条件输入。条件输入不自动保证物理可行性。

去噪扩散隐式模型（Denoising Diffusion Implicit Model，DDIM）构造与上述训练方式兼容的另一类采样过程，可使用较少的噪声步；在其确定性设定下，固定初始噪声和条件便确定了输出，样本多样性来自不同的初始噪声。减少步数仍可能损失质量。潜空间扩散先用自编码器压缩数据，在潜空间训练并采样，最后解码回数据空间，以降低高分辨率任务的计算量。这些变体仍需验证条件一致性；潜空间方法还需检查压缩和解码是否丢失关键细节。

\section{任务选择与验证边界}

VAE、GAN 和扩散模型体现了不同的计算折中。需要连续潜变量和直接的近似后验编码时，VAE 的结构较方便；若用于补全，应把已观测部分与掩码等信息纳入条件或推断过程。需要快速采样、且重视视觉或波形逼真度时，可以考虑 GAN，同时检查训练稳定性和模式坍塌。复杂条件生成也可考虑扩散模型，但多步采样通常带来较高成本；覆盖性和质量仍要依任务验证，不能仅凭模型类别判断。若需要便于计算的显式似然或可逆映射，可进一步了解自回归模型或正规化流。若要做异常检测，还须验证所用分数能否区分目标异常，不能直接把高似然当作正常、低似然当作故障。

工业应用至少要检查四类性质：
\begin{enumerate}
  \item \textbf{真实性：}边缘分布、相关性、频谱或缺陷几何是否接近真实数据。
  \item \textbf{覆盖性：}不同设备、工况和故障类型是否都出现，而不是重复少数模板。
  \item \textbf{条件与物理一致性：}生成结果是否符合输入条件、守恒关系、边界条件和状态转移约束。
  \item \textbf{下游价值：}加入合成数据后，独立真实测试集上的召回、误报和校准是否改善。
\end{enumerate}

约束可以加入生成损失，例如
\begin{equation}
\mathcal L=\mathcal L_{\mathrm{gen}}+\lambda_{\mathrm{phys}}\E[\|r(x,c)\|_2^2],
\end{equation}
这里按最小化损失的约定书写，$\mathcal L_{\mathrm{gen}}$是相应生成模型的训练损失，$\lambda_{\mathrm{phys}}\geq0$；$x$是训练时可计算的生成样本或干净样本估计，$r(x,c)$是其物理残差，期望对所用条件及采样随机性取平均。惩罚过弱可能保留不可行样本，过强则可能损害分布覆盖。软惩罚本身不保证每个样本严格满足约束，应分别报告真实性、多样性和约束违反率。

生成数据还可能记住训练集中的罕见记录，或改变真实故障的先验比例。测试集必须保持真实、独立，并且不能参与生成模型训练或选择。
例如振动窗口来自同一设备或同一次采集时，应先按设备、采集批次或时间段划分训练集、验证集和测试集，再在各集合内切窗，避免高度相似的窗口跨集合泄漏。模型选择和超参数调整使用验证集，测试集用于最终评价。
评价时应与“只用真实数据”的基线比较，匹配训练预算，并在不同随机种子、设备和故障类型上重复。生成样本的数量增加、训练损失下降或单张图片更清晰，都不能单独证明下游收益。

\section*{辅助学习材料}
\begingroup
\emergencystretch=3em
\begin{itemize}
\item \textbf{What are Diffusion Models?}；作者/平台：Lilian Weng；英文技术教程。可对照前向加噪、反向去噪、条件引导和潜空间扩散。\href{https://lilianweng.github.io/posts/2021-07-11-diffusion-models/}{在线阅读}。
\item \textbf{From Autoencoder to Beta-VAE}；作者/平台：Lilian Weng；英文技术教程。可对照 VAE、ELBO 和重参数化。\href{https://lilianweng.github.io/posts/2018-08-12-vae/}{在线阅读}。
\item \textbf{Machine Learning 2021 Spring：GAN}；作者/平台：李宏毅，台湾大学课程网站；提供中文视频入口和讲义。可对照 GAN 的生成器、判别器和训练过程。\href{https://speech.ee.ntu.edu.tw/~hylee/ml/2021-spring.php}{课程页面（GAN 小节）}。
\end{itemize}
\endgroup

\paragraph{本章小结。}
生成任务首先要明确输出分布和条件。VAE 通过潜变量和 ELBO 学习可采样表示，GAN 通过判别器提供生成训练信号，但可能出现训练不稳定和模式坍塌；扩散模型通过逐步去噪生成样本，并可结合条件控制。模型选择不能只看样本外观。真实性、覆盖性、条件与物理一致性，以及独立真实数据上的下游结果共同决定生成模型是否适用。
